\section{Typed Model Foundation and Claim Boundary}
\label{sec:oc133-integrated-scientific-closure}
This chapter states the model-core boundary inside the monograph's main scientific argument. It introduces the typed foundation, proof obligations, finite-check semantics, replay QA rows, comparator claims, and review boundaries as scientific material: definitions first, obligations second, evidence third, and limits fourth.

The public manuscript identifier is OC Core 1.3.3; its Concept DOI is printed on the title page. Operational archive records, repository tags, and platform-specific release events are not part of the scientific argument. They are referenced at the end of the manuscript only where a reader needs the public repository route or citation identity.

\subsection{Evidence Route for the 1.3.3 Model-Core Chapters}
The reader should treat the 1.3.3 chapters as the public scientific route through the proof, finite-model, validation, comparator, and review corpus. The release promotes bounded model-core claims where the artifacts contain explicit evidence. It does not promote complete scientific coverage or unbounded cross-science comparison claims; those obligations remain in the background science program until they are literally evidenced.

\subsection{How Evidence Enters the Model}
The scientific text binds research artifacts, proof registers, Lean and finite evidence, replay QA tables, novelty/comparator rows, and review closures to readable sections. A claim is introduced in prose before the reader is asked to inspect a theorem identifier, replay row, or machine-readable evidence file.

\begin{itemize}[leftmargin=1.8em]
\item \textbf{Manuscript-quality rule:} internal routing memos, raw registers, glued appendices, page-number resets, absent front matter, and absent dedication create scientific reopening conditions
\item \textbf{Corpus coverage gate:} every promoted 1.3.3 science surface is included, cited, or excluded with reason
\item \textbf{Evidence coverage gate:} major theorem, evidence, review, and journal anchors must appear in the manuscript as readable science before machine evidence is consulted
\item \textbf{Editorial adversarial-review gate:} editorial reviewers must return zero critical/high findings before publication replacement
\end{itemize}
\subsection{Scientific Support Obligations}
The 1.3.3 publication standard is deliberately positive. A public scientific document does not pass merely because it avoids forbidden words, stale metadata, or route-sheet fragments. It must also perform its role. For every major claim the manuscript has to say what is being claimed, why the claim matters, which formal or empirical evidence supports it, how that evidence relates to prior art and the current scientific picture, what simulation, finite semantic, or replay route can check it, and what boundary or falsifier would reopen it.

This positive obligation is stricter than a normal release checklist. The manuscript must show that the current science has actually been integrated into the public text: typed foundation, theorem/proof route, Lean subset, finite semantic checks, target-blind numeric rows, negative controls, falsifiers, prior-art comparison, phenomenon coverage, adversarial review, reproducibility, and journal owner-review preparation. If any of those surfaces exists only as a hidden file or a raw register with no reader-facing explanation, the manuscript has not fulfilled its scientific task.

The current science-state register therefore has two jobs. First, it records the internal maturity vector: model foundation, formal evidence, empirical evidence, prior-art alignment, corpus completeness, process visibility, and editorial review state. Second, it constrains external speech. Public wording must derive from the actual research state rather than from ambition, marketing pressure, or local process vocabulary. The reader sees the scientific result; the machine keeps the stronger internal proof that the result is supported, current, and bounded.

This section is included so future OC and non-OC releases cannot repeat the failure pattern that produced a public archive record from metadata, process fragments, or an appended delta. The standard is now part of the manuscript machinery: manuscript integration, editorial review, package construction, public metadata, and post-release verification must all preserve the same positive obligations.

\subsection{Evidence Coverage Map}
For this release the evidence-inclusion rule is applied as a coverage map, not as a cosmetic checklist. Each public scientific surface has a required reader payoff and a required evidence payoff. The model chapter must teach typed carriers, realizations, liveness, death, residue, morphisms, boundaries, operators, cycle modes, dimension, and K-level semantics. The proof chapter must connect theorem names to assumptions, definitions, proof sheets, Lean declarations, finite witnesses, dependency references, and counterexample boundaries. The empirical chapter must distinguish target-blind replay QA and artifact-integrity examples from future full-domain evidence-boundary review.

Prior-art and novelty discussion has its own positive obligations. It must identify overlap with existing traditions, name the residual delta that OC claims under declared assumptions, and refuse unsupported priority or unrestricted comparative claim language. Phenomenon coverage has a parallel obligation: it must state what the model instance explains, what observable or replay route exists, which comparator is relevant, what negative control or falsifier can reopen the claim, and whether the phenomenon is promoted now or left as a research target.

The reviewer chapter therefore cannot be a list of issues. It must show why a hostile objection is serious, which public claim it threatens, which evidence answers it, and what residual risk would reopen it. The methods chapter cannot be a list of commands; it must tell the reader which claim each command checks, which input it consumes, which output or hash it should produce, and what scientific interpretation follows from a mismatch. The journal package map cannot imply submission; it must show preparation readiness and the additional owner/editorial action required before an actual submission.

\begin{itemize}[leftmargin=1.8em]
\item \textbf{Model integration payoff:} the reader can reconstruct the tuple-to-boundary-to-operator route without opening a machine register
\item \textbf{Proof integration payoff:} the reader can move from a promoted theorem name to assumptions, proof sheet, Lean subset, finite witness, and boundary
\item \textbf{Empirical integration payoff:} the reader can see formula, snapshot, split, prediction, observation, uncertainty, residual, comparator, negative control, falsifier, and replay hash where empirical promotion is claimed
\item \textbf{Prior-art payoff:} the reader can distinguish overlap, residual delta, non-novelty boundary, and unsupported priority claim
\item \textbf{Phenomenon payoff:} the reader can tell what is explained, what is only protocol-ready, and what would falsify the explanation
\item \textbf{Editorial payoff:} the reader receives one coherent manuscript voice rather than an old volume plus appended control material
\item \textbf{Release payoff:} public metadata, DOI, archive files, and document surfaces describe the same bounded scientific object
\end{itemize}
These payoffs are deliberately measurable. Missing anchors, absent title/front matter, absent dedication, stale version identity, weak literature synthesis, missing visual pedagogy, unsupported public claims, over-repeated boilerplate, register-dominated prose, non-current review results, and package/public metadata mismatch are all scientific reopening conditions. The public artifact must be pleasant enough to read and strict enough to audit; either failure is a scientific publication failure.


\section{Scientific Reading Protocol}
\label{sec:oc133-reader-contract-scientific-route}
Audience. The primary audience is a mixed external-review group: formal-methods readers, systems-theory readers, domain scientists, journal editors, and technically literate institutional readers. The monograph therefore cannot assume that every reader begins with the same mathematical or domain background.

Purpose. This manuscript explains what OC Core 1.3.3 claims, why the claims are bounded, how the typed model is organized, where proof and executable evidence live, which empirical lanes are replayable, and which broader full-science obligations remain future work. Because the artifact is public science, intelligibility is a release requirement.

Construction. The recommended reader path is orientation, formal model, theorem/proof closure, Lean and finite-model evidence, target-blind empirical rows, prior-art comparison, phenomenon coverage, reviewer objections, reproducibility, and release governance. This order is deliberate: a reader needs the model before the proof register and the proof register before judging replay artifacts.

Didactic rule. Every major section must answer four questions: what is being claimed, why it matters, what evidence supports it, and what would falsify or limit it. If a section only lists identifiers, hashes, or rows, the detailed material belongs in the evidence package or appendix rather than in the main explanatory path.


\section{Literature and Prior-Art Position}
\label{sec:oc133-literature-prior-art-position}
The monograph carries 129 bibliography entries and source-backed comparator anchors with overlap and residual-delta notes. The literature layer locates OC against emergence, phase transitions, autopoiesis, dynamical systems, RAF chemistry, information and complexity measures, identity over time, hybrid systems, formal methods, systems engineering, and reproducibility practice. The release uses literature to bound novelty, not to claim absence of all prior work.

A scientific reviewer should therefore read the comparator material as a map of overlap and residual delta. Where OC integrates or rephrases an existing tradition, the monograph says so. Where a broader priority or unrestricted comparative claim is not evidenced, it stays outside the promoted release surface.


\section{Figure Route and Design Logic}
\label{sec:oc133-figure-route-design-logic}
The full corpus includes 32 public figure entries across the theorem roadmap, worked examples, and figure atlas. The figure layer teaches tuple structure, thresholds, K-level transitions, lifecycle, boundaries, domain examples, and theory/reproducibility movement. Figures are explanatory anchors, not decoration; they are staged so that the main proof path remains readable while visual readers can still follow the structural grammar.

The design rule for 1.3.3 is simple: every public document must begin with a title page, dedication, version, DOI, abstract or reader contract, and table of contents, and every long register must either be explained in prose or moved to an appendix or evidence bundle.


\section{Integrated 1.3.3 Scientific Argument Before Evidence Rows}
\label{sec:oc133-integrated-argument-before-evidence-rows}
Before the monograph names individual evidence rows, it states the integrated argument in ordinary scientific prose. OC Core 1.3.3 promotes a bounded model-core claim: a continuum is treated as a typed object whose carrier, realization, lawful possibility, liveness, residue, boundary behavior, operator behavior, cycle mode, dimension semantics, and K-level position can be inspected under declared assumptions. The release does not ask the reader to trust an inventory. It asks the reader to test whether those typed components solve specific ambiguity classes that appear when persistence, death, rebirth, boundary, update, and cross-level reduction are discussed without a common grammar.

The proof layer then gives the argument its ceiling. The promoted theorem surface is not a claim that OC has completed every possible domain science. It is a claim that named theorem statements have named assumptions, proof sheets, Lean-subset references where available, finite semantic witnesses, negative controls, and counterexample boundaries. A theorem identifier is therefore a citation handle, not the argument itself: the argument is the chain from definition to lemma to proof sheet to finite witness to explicit reopening condition.

The empirical layer has the same discipline. Physics, chemistry, biology, systems, and mathematics rows are target-blind or replay-QA examples with formulas, snapshots, comparators, residuals, negative controls, falsifiers, and hashes. They show how the OC release binds a claim to a replayable evidence route. They do not promote scoped numeric closure, final full-scope synthesis status, or unrestricted comparative claim over modern science. Those broader ambitions remain a background research program until they can be supported at the same or higher evidential standard.

The comparator layer prevents novelty theater. General systems theory, autopoiesis, dynamical systems, category and topos formalisms, RAF theory, complexity and information measures, identity theory, systems engineering, and reproducible-research practice are treated as live comparators. For each tradition, the release first records overlap, then states the residual delta, and finally names the claim that OC must not make. This order is essential: accepted overlap comes before residual contribution.

The row-based sections that follow are therefore not a substitute for the monograph. They are an audit trail. A reader who wants the conceptual argument should read the prose spine first; a reviewer who wants to attack exact wording can then use the rows to locate the proof sheet, finite case, replay row, comparator source, or reopening rule that carries the challenged sentence.
